\documentclass[11pt]{article}
% Source-PDF: 动态规划 DYNAMIC PROGRAMMING/队列优化一类动态规划的优化 - 林涛.pdf
\usepackage{oipl-vn}

\title{Phương pháp tối ưu hóa cho một lớp bài toán quy hoạch động}
\author{Lin Tao}
\date{}

\begin{document}
\maketitle

\origtitle{Optimization Methods for a Class of Dynamic Programming Problems}
\sourcepath{Dynamic Programming / Queue-optimized DP - Lin Tao.pdf}

\section*{Tóm tắt}

Hiện nay các bài toán quy hoạch động đã xuất hiện rất thường xuyên. Phương
trình chuyển trạng thái thường không khó để lập ra, nhưng phần tối ưu thuật
toán lại hay gây trở ngại. Bài viết này thông qua một số bài toán điển hình
trong vài năm gần đây để thảo luận một phương pháp tối ưu cho một lớp bài
toán quy hoạch động. Các tối ưu trong những bài này ít nhiều đều lấy cảm hứng
từ bài \emph{Batch} của IOI 2002; trước khi đọc, có thể tham khảo báo cáo lời
giải của Yu Wei.

\section{Mở đầu}

Có một lớp bài toán quy hoạch động mà biến quyết định có một số tính đơn điệu
nhất định. Dựa vào đặc điểm này, ta có thể tối ưu việc chọn biến quyết định,
từ đó giảm độ phức tạp thời gian của thuật toán. Các ví dụ sau minh họa cách
áp dụng tư tưởng đó.

\section{Euro (Balkan 2003)}

\subsection*{Đề bài}

Cho một dãy $a_1,a_2,\ldots,a_n$ và một số $T$. Cần chia dãy thành một số đoạn.
Mỗi đoạn có một trọng số $F$. Nếu chia $a_i,\ldots,a_j$ thành một đoạn thì
trọng số đoạn này được tính bởi
\[
  F=\operatorname{Sum}[i,j]\cdot j-T,
  \qquad
  \operatorname{Sum}[i,j]=\sum_{k=i}^{j} a_k.
\]
Hãy tìm một cách chia sao cho tổng trọng số của các đoạn là lớn nhất.

Ràng buộc:
\[
  1\le n\le 34567,\qquad
  -1000\le a_i\le 1000\quad (i=1,2,\ldots,n).
\]

\subsection*{Phân tích}

Rất dễ nghĩ đến quy hoạch động. Gọi $F[i]$ là tổng trọng số tối ưu khi chia
$i$ phần tử đầu của dãy, và $t[i]$ là tổng của $i$ phần tử đầu. Khi đó phương
trình chuyển trạng thái là
\[
  F[i]=\max_{0\le k<i}\{F[k]+(t[i]-t[k])\cdot i-T\},
  \qquad 0<i\le n.
\]
Điều kiện biên là $F[0]=0$. Độ phức tạp thời gian của thuật toán trực tiếp là
$O(n^2)$. Liệu có thể tối ưu thêm hay không?

Trước hết, ta phân tích cấu trúc của một phương án chia tối ưu. Giả sử có một
đoạn $a_i,\ldots,a_j$ với $j<n$. Nếu tồn tại $i\le k\le j$ sao cho tổng của
$a_i,\ldots,a_{k-1}$ là âm, còn tổng của $a_k,\ldots,a_j$ là dương, thì cách
chia này chắc chắn không tối ưu. Lý do là ta có thể giữ nguyên các đoạn khác,
tách riêng $a_i,\ldots,a_{k-1}$, đồng thời gộp $a_k,\ldots,a_j$ vào đoạn kế
tiếp. Khi đó phần có tổng âm được nhân với một chỉ số nhỏ hơn, còn phần có
tổng dương được nhân với một chỉ số lớn hơn, nên tổng trọng số tăng lên.

Vì vậy ta có kết luận sau: nếu $a_i,\ldots,a_j$ với $j<n$ là một đoạn con
trong phương án chia tối ưu, thì nhất định phải có
\[
  t[j]\le t[i-1].
\]
Nếu không, ta sẽ tìm được một vị trí $k$ như trên để cải thiện phương án.

Tiếp theo, xét biến quyết định $k$ của $F[i]$. Đặt
\[
  S[k,i]=F[k]+(t[i]-t[k])\cdot i-T.
\]
Với mọi $0\le k_1<k_2<i$, ta có
\[
\begin{aligned}
  S[k_1,i]&=F[k_1]+(t[i]-t[k_1])\cdot i-T,\\
  S[k_2,i]&=F[k_2]+(t[i]-t[k_2])\cdot i-T,
\end{aligned}
\]
nên
\[
  S[k_1,i]-S[k_2,i]
  =F[k_1]-F[k_2]+(t[k_2]-t[k_1])\cdot i.
\]

\paragraph{Trường hợp 1: $t[k_2]\ge t[k_1]$.}
Khi đó tổng của $a_{k_1+1},\ldots,a_{k_2}$ là dương. Do đó, dù dãy
$a_1,\ldots,a_{k_2}$ được chia như thế nào, trong phương án tương ứng tất yếu
xảy ra một trong hai tình huống: có một đoạn có tổng dương, hoặc phần đuôi
của một đoạn có tổng dương. Cả hai tình huống đều cho thấy đây không phải
phương án tối ưu. Nói cách khác, $k_2$ không có ý nghĩa làm quyết định.

\paragraph{Trường hợp 2: $t[k_2]<t[k_1]$.}
Nếu
\[
  S[k_1,i]-S[k_2,i]
  =F[k_1]-F[k_2]+(t[k_2]-t[k_1])\cdot i>0,
\]
tức là
\[
  i<\frac{F[k_2]-F[k_1]}{t[k_1]-t[k_2]},
\]
thì với mọi $j\le i$ ta cũng có
\[
  S[k_1,j]-S[k_2,j]>0.
\]
Như vậy biến quyết định có tính đơn điệu.

Đặt
\[
  g(j,k)=\frac{F[k]-F[j]}{t[j]-t[k]}.
\]
Với $0<k_1<k_2<k_3<\cdots<k_m$, ta có thể duy trì một hàng đợi sao cho
\[
  t[k_1]>t[k_2]>t[k_3]>\cdots>t[k_m],
\]
đồng thời bảo đảm
\[
  i<g(k_2,k_1)<g(k_3,k_2)<g(k_4,k_3)<\cdots<g(k_m,k_{m-1}).
\]
Khi đó
\[
  S[k_1,i]>S[k_2,i]>S[k_3,i]>\cdots>S[k_m,i],
\]
và quyết định tối ưu của $F[i]$ là $k_1$:
\[
  F[i]=F[k_1]+(t[i]-t[k_1])\cdot i-T.
\]

\subsection*{Cài đặt}

Các thao tác duy trì hàng đợi như sau.

\begin{enumerate}
  \item Thêm một phần tử mới $k$. Vì $k>k_m$, ta đặt $k$ vào cuối hàng đợi.
  Nếu
  \[
    g(k,k_m)<g(k_m,k_{m-1}),
  \]
  thì xóa $k_m$, tức $m\leftarrow m-1$, và lặp lại cho đến khi
  \[
    g(k_m,k_{m-1})<g(k,k_m).
  \]

  \item Xóa phần tử đầu. Nếu
  \[
    g(k_2,k_1)<i,
  \]
  thì xóa $k_1$ và lặp lại cho đến khi
  \[
    g(k_2,k_1)>i.
  \]
\end{enumerate}

Vì tổng số phần tử được thêm và bị xóa đều không vượt quá $n$, độ phức tạp
thời gian của thuật toán là $O(n)$.

\section{Giá trị trung bình lớn nhất (USACO)}

\subsection*{Đề bài}

Cho một dãy $a_1,a_2,\ldots,a_n$. Hãy tìm một đoạn có độ dài ít nhất là $L$
sao cho giá trị trung bình của đoạn đó là lớn nhất.

\subsection*{Phân tích}

Giả sử trong các đoạn kết thúc tại $a_i$, đoạn có giá trị trung bình lớn nhất
bắt đầu từ phần tử thứ $F[i]$. Khi đã biết $F[i]$, ta xét cách tìm $F[i+1]$.

Giả sử $F[i+1]<F[i]$. Điều đó có nghĩa là giá trị trung bình của đoạn từ
$F[i+1]$ đến $F[i]-1$ lớn hơn giá trị trung bình của đoạn từ $F[i]$ đến
$i+1$. Tuy nhiên, theo định nghĩa, giá trị trung bình của đoạn từ $F[i+1]$ đến
$F[i]-1$ không thể lớn hơn giá trị trung bình của đoạn từ $F[i]$ đến $i$.
Vì vậy đoạn tốt nhất kết thúc tại $i+1$ sẽ không vượt qua đoạn tốt nhất kết
thúc tại $i$, nên trường hợp này không có ý nghĩa.

Từ kết luận trên, khi tính giá trị trung bình lớn nhất, ta chỉ cần quan tâm
tới các đoạn kết thúc tại $i+1$ trong trường hợp $F[i+1]\ge F[i]$.

Giả sử ta đã tính được $F[i]$. Trong quá trình tính các trạng thái về sau,
vì vị trí bắt đầu của đoạn liên tục tăng, ta cần xét khi nào nên tăng và tăng
bao nhiêu. Gọi giá trị trung bình lớn nhất hiện tại là $x$, vị trí bắt đầu
hiện tại là $p$, và vị trí cuối là $i$. Ta xét đoạn các vị trí có thể xóa từ
$p$ đến $i-L$. Nếu giá trị trung bình từ $p$ đến một vị trí nào đó nhỏ hơn
giá trị trung bình lớn nhất hiện tại, thì phần đó không còn ý nghĩa và có thể
xóa trực tiếp.

Giả sử giá trị trung bình của các đoạn
\[
  p\ldots p,\quad p\ldots p+1,\quad p\ldots p+2,\ldots
\]
lần lượt là $x_0,x_1,x_2,\ldots$. Nếu $x_i>x_{i+1}$, thì sau khi xóa đoạn
$p$ đến $p+i$, chắc chắn ta cũng sẽ xóa tiếp $p+i+1$. Sau khi xóa $p+i+1$,
việc có nên xóa $p+i+2$ hay không lại cần được tính lại.

Nhưng việc tính lại này là không cần thiết, vì ta có thể tránh trước. Nguyên
nhân phải tính lại là trong dãy xuất hiện hiện tượng $x_i>x_{i+1}$. Do đó ta
duy trì một dãy vị trí
\[
  P_1,P_2,P_3,\ldots,
\]
trong đó $x_i$ là giá trị trung bình của đoạn từ $P_i$ đến $P_{i+1}-1$, và
các giá trị này thỏa
\[
  x_1<x_2<x_3<x_4<\cdots.
\]
Vì phần có thể xóa không được vượt quá $i-L$, nên khi $i$ tăng, phần có thể
xóa cũng tăng. Mỗi khi dãy này được mở rộng, ta phải duy trì tính chất của
dãy $P$.

Giả sử ban đầu ta có $P_1,P_2,P_3,\ldots,P_m$, và vị trí mới được phép xóa
tăng đến $P_{m+1}-1$. Nếu $x_m>x_{m-1}$, tính chất của dãy không đổi nên
không cần xử lý. Nếu $x_m<x_{m-1}$, ta gộp đoạn từ $P_m$ đến $P_{m+1}-1$ vào
đoạn trước đó, xóa $P_m$, rồi lặp lại thao tác này.

Bản chất của thao tác trên là: nếu giá trị trung bình của đoạn cuối nhỏ hơn
giá trị trung bình của đoạn áp chót, ta gộp hai đoạn đó thành một đoạn. Vì
mỗi phần tử trong dãy chỉ được thêm hoặc xóa một lần, độ phức tạp của thuật
toán là $O(n)$.

Bạn đọc cũng có thể tham khảo phương pháp kết hợp hình học và đại số mà Zhou
Yuan từng nhắc đến ở trại mùa đông để giải bài này.

\section{Máy tính kiểu mới (OIBH)}

\subsection*{Đề bài}

Cho một dãy số. Ta cần thay đổi dãy sao cho nó thỏa điều kiện sau: bắt đầu từ
số đầu tiên, đọc vào một số $i$, sau đó đọc ra đúng $i$ số đứng ngay sau nó;
tiếp tục đọc vào một số $j$, rồi đọc ra đúng $j$ số đứng ngay sau nó; và cứ
lặp lại như vậy, sao cho toàn bộ dãy vừa được đọc hết đúng lúc. Chi phí thay
đổi dãy là tổng các giá trị tuyệt đối của chênh lệch trước và sau khi thay
đổi trên mọi phần tử. Hãy tìm một phương án có chi phí thay đổi nhỏ nhất.

\subsection*{Phân tích}

Khi đọc đề lần đầu, ta rất dễ nghĩ ra phương trình quy hoạch động. Gọi $F[i]$
là chi phí nhỏ nhất để đoạn bắt đầu tại vị trí $i$ được đọc hết đúng cách.
Khi đó
\[
  F[i]=\min_{k>i}\{F[k]+|a_i-(k-i-1)|\}.
\]
Độ phức tạp của thuật toán này hiển nhiên là $O(n^2)$.

Đặt
\[
  S[i,k]=F[k]+|a_i+i+1-k|.
\]
Nếu với $k<j$ có
\[
  S[i,k]<S[i,j],
\]
thì
\[
  |a_i+i+1-k|-|a_i+i+1-j|<F[j]-F[k].
\]

Tập nghiệm của bất đẳng thức này có thể có dạng
\[
  a_i+i+1<x.
\]
Dựa vào tính đơn điệu này, ta thử xây dựng một dãy
\[
  k_1>k_2>k_3>\cdots
\]
sao cho tập nghiệm của bất đẳng thức
\[
  S[i,k_{r+1}]<S[i,k_r]
\]
có dạng
\[
  a_i+i+1<x_r,
\]
và đồng thời
\[
  x_1>x_2>x_3>x_4>\cdots.
\]
Khi đã có một dãy như vậy, để tính $F[i]$ ta không cần liệt kê mọi biến quyết
định $k$. Thay vào đó, ta nhị phân tìm chỉ số lớn nhất $r$ sao cho
\[
  a_i+i+1<x_r,
\]
rồi lấy $k_r$ làm quyết định.

Dĩ nhiên, dãy $k_1,k_2,k_3,\ldots$ được tăng dần trong quá trình xử lý. Khi
thêm một phần tử mới $k_{m+1}$, ta giải bất đẳng thức
\[
  S[i,k_{m+1}]<S[i,k_m]
\]
và thu được tập nghiệm $a_i+i+1<x_{m+1}$. Nếu $x_{m+1}<x_m$ thì tính chất của
dãy không đổi; ngược lại, ta xóa $k_m$ và lặp lại thao tác này cho đến khi
dãy lại thỏa
\[
  x_1>x_2>x_3>x_4>\cdots.
\]

Độ phức tạp của các thao tác duy trì là $O(n)$, nên độ phức tạp toàn bài là
$O(n\log n)$.

Thực ra bài này còn có một cách giải tốt hơn bằng đồ thị, có thể giảm độ
phức tạp xuống $O(n)$.

\section{Tổng kết}

Với các bài toán trên, việc viết phương trình quy hoạch động không khó. Tuy
nhiên, theo cách thông thường, biến quyết định cần được liệt kê. Bài viết này
tập trung vào một số tính đơn điệu của biến quyết định, rồi tối ưu thêm ở
khâu chọn quyết định, từ đó đạt mục tiêu giảm độ phức tạp thời gian.

Các lời giải trên có một số điểm chung, cũng là các điểm chung của kiểu tối
ưu này:
\begin{enumerate}
  \item Từ phương trình quy hoạch động, phát hiện một số tính đơn điệu của
  biến quyết định.
  \item Dùng tính đơn điệu đó để duy trì một hàng đợi thể hiện chính tính
  chất đơn điệu.
  \item Dùng hàng đợi này để chọn quyết định.
\end{enumerate}

\end{document}
